%% ============================================================
%% 片段：s03b-parabolic   对应原文 PDF 第 48--61 页
%% 本文件由 gw/convert.py 从 arXiv 源码（hep-th/0612073v2）自动生成：
%%   * 公式/编号/引用/图 已经锁定为原文内容，请勿改动任何数学内容与编号；
%%   * 只需把英文正文译为中文，并在每节开头加“本节重点”、关键处加译注。
%% ============================================================


\subsection{抛物丛}

%% ==== 原文 §3.4：PDF 第 48 页（印刷第 47 页）====
\label{sec:3.4}

至此，我们或许应当解释一个在数学理论中通常被当作出发点、但迄今被我们隐藏起来的概念。这就是抛物丛（parabolic bundle）的概念，它与抛物希格斯丛（parabolic Higgs bundle）不同。所有陈述在多个点处带有抛物结构的情形都有显然的类比，但为简单起见，我们主要考虑一个点的情形。在另行说明之前，我们只考虑规范场论，不含希格斯场。

作为动机，我们回到第~\ref{sec:3.3}~节的 $SU(2)$ 例子。丛 ${\rm ad}(E)$ 的一个全纯截面 $s$ 是一个被 $\bar\partial_A$ 算符零化的 ${\rm ad}(E)$-值函数。给定该算符的具体形式~\eqref{eq:3.20}，这意味着在 $z=0$ 附近，
% 式 (3.33)
\setcounter{equation}{32}
\begin{equation}\label{eq:3.33}
  s=\left(\begin{matrix}u & v(\bar z z)^{y}\\ w(\bar z z)^{-y} & -u\end{matrix}\right),
\end{equation}
其中 $u,v,w$ 是普通全纯函数。若希望 $|s|$ 在 $z\to 0$ 时有界，则要求 $w(0)=0$。于是在 $z=0$ 处，$\tilde s=f^{-1}sf$ 取形式
% 式 (3.34)
\setcounter{equation}{33}
\begin{equation}\label{eq:3.34}
  \tilde s=\left(\begin{matrix}* & * \\ 0 & * \end{matrix}\right).
\end{equation}

{}从全纯的观点看，实际发生的是：$\alpha$ 的选取在 $p$ 处 ${\ad}(E)$ 的纤维中决定了一个博雷尔子代数（Borel subalgebra）$\mathfrak b$，它由满足 $-i[\alpha,\psi]=\lambda\psi$、$\lambda \geq 0$ 的向量 $\psi\in \mathfrak g$ 张成。式~\eqref{eq:3.34} 表明 $\tilde s(p)$ 取值于 $\mathfrak b$。在这一形式下，该结果对任意 $G$ 成立：若 $\alpha$ 是正则的，则 $\tilde s(p)$ 取值于由 $\alpha$ 决定的博雷尔子群（Borel subgroup）。更一般地，若选取列维子群（Levi subgroup）$\mathbb{L}$ 且 $\alpha$ 是 $\mathbb{L}$-正则的，则由同类的论证，$\tilde s(p)$ 取值于由对 $(\mathbb{L},\alpha)$ 决定的抛物子群（parabolic subgroup）$\EUP$ 的李代数 $\mathfrak p$ 中。

\begin{yizhu}[原文疑误]原文在同一段先后用两种说法：先引入由 $\alpha$ 决定的博雷尔子代数 $\mathfrak b$，并用式~\eqref{eq:3.34} 说明 $\tilde s(p)$ 取值于 $\mathfrak b$；随后却说 $\tilde s(p)$ 取值于由 $\alpha$ 决定的博雷尔子群（Borel subgroup）。$\tilde s$ 是伴随丛 ${\rm ad}(E)$ 的截面，$\tilde s(p)$ 是李代数元素，而式~\eqref{eq:3.34} 显示的也确实是一个矩阵，故后一处应为博雷尔子代数。原文如此（PDF 第 49 页，印刷第 48 页），译文照录原说法。\end{yizhu}

$G$-丛 $E\to C$ 在一点 $p$ 处的抛物结构（parabolic structure）选取，无非是把 $E$ 在 $p$ 处的结构群约化到抛物子群 $\EUP$。我们刚刚看到的是：对 $\mathbb{L}$-正则的 $\alpha$，一个在微分几何中由 $p\in C$ 附近的 $A=\alpha\,d\theta+\dots$ 描述的带奇点丛，在复几何中对应于一个在 $p$ 处带有抛物结构选取的丛。

Mehta 与 Seshadri 的一个定理~\cite{ref50}（它推广了无抛物结构情形下 Narashimhan 与 Seshadri 的定理~\cite{ref51}）把这一点纳入系统的框架。该定理在稳定抛物 $G_\C$-丛与带有熟悉的奇点 $A=\alpha\,d\theta+\dots$ 的平坦 $G$-丛之间建立一一对应。在对应的左侧，考虑一个全纯 $G_\C$-丛，其在某点 $p\in C$ 处的结构群被约化到抛物子群 $\EUP$。令 $\mathbb{L}$ 为 $\EUP$ 的列维子群，并选取 $\alpha$，使得 $\EUP$ 按通常方式由对 $(\mathbb{L},\alpha)$ 决定。我们假定 $\alpha$ 足够一般，使得 $G$ 中与 $U=\exp(-2\pi\alpha)$ 交换的子群恰为 $\mathbb{L}$。我们称这样的 $\alpha$ 是严格 $\mathbb{L}$-正则的。对每个这样的 $\alpha$，带有 $\EUP$ 型抛物结构的丛都有自然的稳定性（stability）概念。\footnote{此外，在等价意义下，该概念在把 $\alpha$ 平移一个格向量时不变。这里所说的等价涉及对丛作赫克修改（Hecke modification）。}

\begin{yizhu}[原文疑误]原文把 Narasimhan 拼作“Narashimhan”（“Narashimhan 与 Seshadri 的定理”）。正确拼写为 Narasimhan——即与 Seshadri 共同证明稳定丛模空间基本定理的数学家；本条定理在文献中通常称为 Narasimhan–Seshadri 定理。译文照录原拼写。原文如此（PDF 第 49 页，印刷第 48 页）。\end{yizhu}

\begin{yizhu}[易错点]原文只说“对每个严格 $\mathbb{L}$-正则的 $\alpha$ 存在自然的稳定性概念”，却未给出定义。标准定义是抛物稳定性的斜率判据：$\EUP$ 的幂幺根给出 $E$ 在 $p$ 处纤维上的滤过，$\alpha$ 给出相应的抛物权重 $\alpha_1>\alpha_2>\dots$；对每个保持该滤过的子丛 $F\subset E$，要求其抛物斜率（$\deg F$ 加上按 $\alpha_i$ 加权的各级商丛秩之和）小于 $E$ 的抛物斜率（此处因 $\deg E=0$ 即为 $0$）。取 $\EUP$ 为博雷尔子群时即退回通常的稳定性；$\alpha$ 的格平移不变性则由脚注提到的赫克修改实现。详见 \cite{ref50}。\end{yizhu}

我们在此不描述这一稳定性条件，尽管它在数学理论中是根本性的。在对应的另一侧，考虑 $C\backslash p$ 上的平坦 $G$-丛，其绕 $p$ 的单值性（monodromy）与 $U=\exp(-2\pi\alpha)$ 共轭。等价地，考虑熟悉的方程
% 式 (3.35)
\setcounter{equation}{34}
\begin{equation}\label{eq:3.35}
  F=2\pi\alpha\delta_p
\end{equation}
的解，模去在点 $p$ 处取值于 $\mathbb{L}$ 的规范变换。我们把这类带单值性的平坦丛的模空间记为 $\M(\alpha;p)$。Mehta 与 Seshadri 的定理断言：稳定抛物 $G_\C$-丛的模空间与 $\M(\alpha;p)$ 相同。无抛物结构情形下的类似定理~\cite{ref51} 则说：$C$ 上稳定 $G_\C$-丛的模空间 $\M$ 与平坦 $G$-丛的模空间相同。

从规范场空间对规范变换群取辛商（symplectic quotient）的观点看，这两个结果都是自然的~\cite{ref52}。联络空间（或带有奇点 $A=\alpha\,d\theta+\dots$ 的联络空间）是辛流形，其辛形式为
% 式 (3.36)
\setcounter{equation}{35}
\begin{equation}\label{eq:3.36}
  \omega=-{1\over
4\pi}\int_C\,\Tr\,\delta A\wedge \delta A
\end{equation}
而矩映射（moment map）为
% 式 (3.37)
\setcounter{equation}{36}
\begin{equation}\label{eq:3.37}
  \mu=-{1\over 2\pi}\int_C\,\Tr\,\epsilon F.
\end{equation}
用适当的规范变换群对联络空间（或带奇点的联络空间）取辛商，得到的是 $\M$（或 $\M(\alpha;p)$）。把该辛商重新解释为对复值规范变换群 $\EUG_\C$ 的商，则这些空间又可解释为稳定丛、或带抛物结构的稳定丛的模空间。其论证与我们在第~\ref{sec:3.1}~节、第~\ref{sec:3.2}~节对希格斯丛详细描述过的相似。

\begin{yizhu}[物理图像补充]式~\eqref{eq:3.36}、\eqref{eq:3.37} 给出同一套辛商的两种描述：矩映射 $\mu=-{1\over 2\pi}\int_C\Tr\epsilon F$ 的零点条件正是 $F=0$，即平坦；对带奇点的联络，零点条件即式~\eqref{eq:3.35} 的 $F=2\pi\alpha\delta_p$。对规范群取辛商得 $\M(\alpha;p)$，改对复化群 $\EUG_\C$ 取商则得稳定抛物丛的模空间（Kempf--Ness）。$\alpha$ 在此同时是奇点强度与 $p$ 处的抛物权重，因此 $\alpha$ 的 $\mathbb{L}$-正则性与稳定性的门槛是同一件事的两面。\end{yizhu}

用 $\M(\alpha;p)$ 记 $C\backslash p$ 上绕 $p$ 有单值性 $U=\exp(-2\pi\alpha)$ 的平坦丛的模空间，这一记号略有误导性，因为 $\M(\alpha;p)$ 并不随 $\alpha$ 光滑变化。其维数依赖于 $G$ 中与 $U$ 交换的子群。只有当 $\alpha$ 被限制为对某个固定的 $\mathbb{L}$ 严格 $\mathbb{L}$-正则时，$\M(\alpha;p)$ 才随 $\alpha$ 光滑变化。为强调这一点，我们把 $\mathbb{L}$ 视为定义的一部分，并将该空间记为 $\M_\mathbb{L}(\alpha;p)$。对给定的 $\mathbb{L}$，严格 $\mathbb{L}$-正则的 $\alpha$ 的空间具有互不相同的连通分支（如我们在第~\ref{sec:2.6}~节所知，它们一般与不同构的抛物子群相联系）。当我们不显式写出 $\mathbb{L}$ 时，意即取 $\mathbb{L}=\TT$。

\begin{yizhu}[易错点]$\M_\mathbb{L}(\alpha;p)$ 的下标 $\mathbb{L}$ 指列维子群，不是 Langlands 对偶群，也不是希格斯侧的下标 $H$。原文随后（式~\eqref{eq:3.38} 一带）指出 $\mathbb{L}=G\cap\EUP$ 同时是 $G$ 的列维子群与 $\EUP$ 的极大紧子群，且 $G_\C/\EUP=G/\mathbb{L}$；后文正是靠这一点把复纤维化~\eqref{eq:3.38} 换成辛纤维化~\eqref{eq:3.39}。第~\ref{sec:3.6}~节、第~\ref{sec:3.7}~节讨论的 $\MH$ 中的 $H$ 指 Hitchin 系统，与 $\mathbb{L}$ 无关，两者通过 $\alpha$ 的极限相联系。\end{yizhu}

\bigskip\noindent{\it 抛物希格斯丛}

现在我们纳入希格斯场，并考虑希格斯丛的类似概念。

在数学理论中，抛物希格斯丛（parabolic Higgs bundle）$(E,\varphi)$ 的概念通常定义如下。$E$ 是上述意义上的抛物丛，其在点 $p$ 处的结构群被约化到 $\EUP$。而微分 $\varphi$ 在 $p$ 处有极点，其留数要求取值于相应的李代数 $\mathfrak p$；此外，该极点的留数与 $\sigma=\half(\beta+i\gamma)$ 具有相同的“特征值”，事实上它等于 $\sigma$ 模 $\mathfrak n$。（其具体例证见式~\eqref{eq:3.25}。）

若 $\alpha$ 与 $\beta+i\gamma$ 是正则的，则无须显式引入抛物结构的概念。对通有的 $(\alpha,\beta,\gamma)$，只需给出丛 $E$ 连同微分 $\varphi$，并要求 $\varphi$ 在 $p$ 处有极点、其留数共轭于 $\half(\beta+i\gamma)$。此时 $\alpha$ 由 $\varphi$ 确定，因为它必须与 $\beta$ 和 $\gamma$ 交换，且其共轭类已知。\footnote{外尔群的作用不带来歧义，因为它在三元组 $(\alpha,\beta,\gamma)$ 上的作用是对角的。}于是抛物结构被确定；而且对通有的参数，我们可以在从不解释丛“具有抛物结构”意味着什么的情况下定义抛物希格斯丛。这实质上就是我们在本课题最初进路中所做的。

\subsection{\texorpdfstring{$\M_\mathbb{L}(\alpha;p)$ 的拓扑}{M_L(α;p) 的拓扑}}

%% ==== 原文 §3.5：PDF 第 51 页（印刷第 50 页）====
\label{sec:3.5}

接下来我们利用上述部分想法，对 $\M_\mathbb{L}(\alpha;p)$ 的拓扑获得粗略的理解。我们旨在为从未接触过这些空间的读者提供关于其拓扑的初步导引，同时也为第~\ref{sec:3.6}~节与第~\ref{sec:3.7}~节所述关于 $\MH$ 的某些结果准备必要的背景。但愿我们所给出的很不完整的解释足以服务这些特定目标。

设给定一个 $G_\C$-丛 $E\to C$ 与 $G_\C$ 的一个抛物子群 $\EUP$，我们希望在某点 $p\in C$ 处选取 $\EUP$ 型抛物结构。我们须在 $p$ 处选取 $G_\C$ 中一个与 $\EUP$ 共轭的子群；所有这类子群构成的空间同构于 $G_\C/\EUP$。这提示 $\M_\mathbb{L}(\alpha;p)$ 应当是底为 $\M$（无抛物结构的稳定 $G$-丛的模空间）、纤维为 $G_\C/\EUP$ 的纤维丛：
% 式 (3.38)
\setcounter{equation}{37}
\begin{equation}\label{eq:3.38}
  \begin{matrix}G_\C/\EUP&\to&
\M_\mathbb{L}(\alpha;p)\\
                                &   & \downarrow \\
                                & & ~\M. \end{matrix}
\end{equation}

在其成立的范围之内（我们马上讨论这一点），该纤维化阐明了 $\M_\mathbb{L}(\alpha;p)$ 的复结构：纤维与底都是复流形，且纤维化是全纯的。

此外，$\M_\mathbb{L}(\alpha;p)$ 具有自然的辛结构，把它解释为带单值性的平坦 $G$-丛的模空间便可方便地理解这一点。其辛形式为 $\omega=-{1\over 4\pi}\int_C\Tr\,\delta A\wedge\delta A$。$\M(\alpha;p)$ 的复结构与辛结构合起来构成 Kähler 结构。为看出 $\M(\alpha;p)$ 的辛结构，\eqref{eq:3.38} 的一个变体更有帮助。复李群理论的一个基本事实是：商 $G_\C/\EUP$ 与 $G/\mathbb{L}$ 相同，其中 $\mathbb{L}=G\cap \EUP$ 是 $G$ 的列维子群，也是 $\EUP$ 的极大紧子群。因此我们可以不用~\eqref{eq:3.38}，而把 $\M_\mathbb{L}(\alpha;p)$ 表示为辛流形的纤维化：
% 式 (3.39)
\setcounter{equation}{38}
\begin{equation}\label{eq:3.39}
  \begin{matrix}G/\mathbb{L}&\to& \M_\mathbb{L}(\alpha;p)\\
                                &   & \downarrow \\
                                & & ~\M. \end{matrix}
\end{equation}
这里底与纤维都是辛流形，下文将用该纤维化描述 $\M_\mathbb{L}(\alpha;p)$ 的辛结构。

\begin{yizhu}[易错点]上文断言 $G_\C/\EUP$ 与 $G/\mathbb{L}$ 相同，这只是在实流形（乃至辛流形）层面的等同：$\EUP$ 是复群，$\mathbb{L}$ 是紧群，$G_\C/\EUP$ 带复结构而 $G/\mathbb{L}$ 只是实流形，二者实维数相同而“复维数”无从比较。因此式~\eqref{eq:3.38} 是复解析描述、式~\eqref{eq:3.39} 是辛描述，二者给出同一个空间；不要把 $G/\mathbb{L}$ 当成复流形去谈全纯纤维化。\end{yizhu}

事实上，纤维化~\eqref{eq:3.38} 与~\eqref{eq:3.39} 恰在我们能假定 $\M$ 中每一点都由稳定（而不只是半稳定）丛代表的程度内成立。否则，可能出现稳定抛物丛 $(E,\EUP)$ 而其底丛 $E$ 并不稳定的情形。此时纤维化失效，因为不能通过先取丛 $E$、再赋予它所有可能的抛物结构的方式来构造稳定抛物丛的模空间 $\M_\mathbb{L}(\alpha;p)$。在 $\M$ 的奇点之外，纤维化~\eqref{eq:3.38} 与~\eqref{eq:3.39} 确实成立，并给出 $\M_\mathbb{L}(\alpha;p)$ 拓扑的良好初步近似。

有一个重要且被广泛研究的情形，其中所有半稳定丛都是稳定的，因此纤维化~\eqref{eq:3.41} 与~\eqref{eq:3.42} 精确成立。这发生在 $G=PSU(N)$ 且 $E$ 在 $H^2(C;\Z_N)$ 中的特征类为 $N$ 阶的情形。一般地，纤维化失效处的余维随 $C$ 的亏格 $g_C$ 或 $G$ 的秩 $r$ 的增大而增大。例如，若 $g_C>2$，或 $g_C=2$、$r>1$，则实余维超过 2。

\begin{yizhu}[原文疑误]此处“纤维化~\eqref{eq:3.41} 与~\eqref{eq:3.42}”应是~\eqref{eq:3.38} 与~\eqref{eq:3.39}（或~\eqref{eq:3.38}、\eqref{eq:3.39} 的辛版本）：\eqref{eq:3.41}、\eqref{eq:3.42} 是上同调正合列，并非纤维化，也与“半稳定丛皆稳定”这一条件无关。原文如此（PDF 第 52 页），照录并加注。\end{yizhu}

若奇点不起重要作用，则可用~\eqref{eq:3.38} 或~\eqref{eq:3.39} 描述 $\M_\mathbb{L}(\alpha;p)$ 的第二上同调群。（我们这样做，是因为它最终有助于理解 $\MH$ 的辛结构，并具体说明仿射外尔群在其上同调上的作用。）我们将主要在假定 $G$ 单连通的情况下进行。此时，$\M_\mathbb{L}(\alpha;p)$ 的上同调的 Leray 谱序列（Leray spectral sequence）起始于
% 式 (3.40)
\setcounter{equation}{39}
\begin{equation}\label{eq:3.40}
  \begin{array}{c|cccc} 2 & H^2(G/\mathbb{L};\mathbb{Z}) & 0 & * & *\\ 1 & 0 & 0 & 0 & 0\\ 0 & \mathbb{Z} & 0 & H^2(\mathcal{M};\mathbb{Z}) & *\\ \hline   & 0 & 1 & 2 & 3\end{array}
\end{equation}

（我们画出了 $\M$ 的第 $q$ 上同调、取值于 $G/\TT$ 的第 $p$ 上同调，$q=0,1,2,3$ 水平排列，$p=0,1,2$ 竖直排列。标有 $*$ 的群的具体形式并不重要。）

\begin{yizhu}[易错点]注意表中指标的排布方向与标准记法相反。按习惯 $E_2^{p,q}=H^p(B;H^q(F))$（$F\to X\to B$ 为纤维化），$p$ 是底 $B=\M$ 的次数、$q$ 是纤维 $G/\TT$ 的次数；此处作者把纤维次数 $p$ 放在竖直方向、底次数 $q$ 放在水平方向（$E_2^{p,q}$ 的行标为纤维、列标为底）。总次数 $p+q=2$ 的格子是 $(p,q)=(2,0)$、$(1,1)$、$(0,2)$，即表中的 $H^2(G/\mathbb{L};\Z)$、$0$、$H^2(\M;\Z)$ 三格；正合列~\eqref{eq:3.41} 正是 $d_2$ 消失后由这三格给出的短正合列。核对时若把行标当成底次数，会错误地以为 $H^2(\M;\Z)$ 与 $H^2(G/\mathbb{L};\Z)$ 的位置颠倒。\end{yizhu}

在二维情形，正如我们将在后文解释的，若 $G$ 单连通，则谱序列中的微分消失，于是该纤维化的谱序列归结为正合列
% 式 (3.41)
\setcounter{equation}{40}
\begin{equation}\label{eq:3.41}
  0\to H^2(\M;\Z)\to
H^2(\M_\mathbb{L}(\alpha;p);\Z)\to H^2(G/\mathbb{L};\Z)\to 0.
\end{equation}
此外，对单连通的 $G$ 有 $H^2(\M;\Z)\cong \Z$，它由某线丛 $\mathfrak L$ 的第一陈类生成~\cite{ref53,ref54}；我们将宽松地称 $\mathfrak L$ 为行列式线丛（determinant line bundle）（对 $G=SU(N)$，$\mathfrak L$ 可定义为某狄拉克算符的行列式线丛）。于是正合列成为
% 式 (3.42)
\setcounter{equation}{41}
\begin{equation}\label{eq:3.42}
  0\to \Z\to H^2(\M_\mathbb{L}(\alpha;p);\Z)\to H^2(G/\mathbb{L};\Z)\to 0.
\end{equation}

\begin{yizhu}[易错点]正合列~\eqref{eq:3.41}、\eqref{eq:3.42} 的成立隐含两点：(i) 该纤维化确为纤维丛（作者在正文中已承认只在 $\M$ 的奇点之外成立）；(ii) 谱序列中所有相关微分 $d_2,d_3$ 消失。原文说“在二维情形……微分消失”，其真正理由在稍后“仿射权格”一段才给出：$H^2(\M(\alpha;p);\Z)$ 含有仿射权格，而若 $d_2$ 非零，该格不可能嵌入。另需注意 $H^2(\M;\Z)\cong\Z$ 与 $H^2(G/\mathbb{L};\Z)$ 均无挠，因此这里的短正合列与“$H^2(X;\Z)$ 的秩及挠部”两种算法一致；对非单连通的 $G$，$\Z$ 与 $\Lambda_{\rm wt}$ 都须换成相应的子格（见~\eqref{eq:3.55} 及其脚注）。\end{yizhu}

对单连通的 $G$，纤维化
% 式 (3.43)
\setcounter{equation}{42}
\begin{equation}\label{eq:3.43}
  \begin{array}{ccc} \mathbb{L} & \to & G\\ & & \downarrow\\ & & G/\mathbb{L}\end{array}
\end{equation}
的谱序列给出 $H^2(G/\mathbb{L};\Z)=H^1(\mathbb{L};\Z)$。对 $\mathbb{L}=\TT$，这给出
% 式 (3.44)
\setcounter{equation}{43}
\begin{equation}\label{eq:3.44}
  H^2(G/\TT;\Z)=H^1(\TT;\Z)=\Lambda_{\rm wt},
\end{equation}
其中 $\Lambda_{\rm wt}$ 是 $G$ 的权格（weight lattice）。这一特定结果（连同同一个格 $\Lambda_{\rm wt}$）无论 $G$ 是否单连通都成立，因为若用 $G$ 的有限覆盖替代 $G$，$G/\TT$ 不变（该覆盖以同样方式扩张 $\TT$，并在商 $G/\TT$ 中消去）。

借助~\eqref{eq:3.44}，对 $\mathbb{L}=\TT$，\eqref{eq:3.42} 等价于
% 式 (3.45)
\setcounter{equation}{44}
\begin{equation}\label{eq:3.45}
  0\to\Z\to
H^2(\M(\alpha;p);\Z)\to \Lambda_{\rm wt}\to 0.
\end{equation}

对更一般的 $\mathbb{L}$，类似的论证给出的是
% 式 (3.46)
\setcounter{equation}{45}
\begin{equation}\label{eq:3.46}
  0\to\Z\to
H^2(\M_\mathbb{L}(\alpha;p);\Z)\to \Lambda_{{\rm wt},\mathbb{L}}\to 0,
\end{equation}
其中 $\Lambda_{{\rm wt},\mathbb{L}}$ 是 $\Lambda_\wt$ 中在 $\mathbb{L}$ 的外尔群作用下不变的子格。

如同任何格的短正合列，\eqref{eq:3.45} 可以分裂，给出
% 式 (3.47)
\setcounter{equation}{46}
\begin{equation}\label{eq:3.47}
  H^2(\M(\alpha;p);\Z)=\Z\oplus \Lambda_{\rm wt}.
\end{equation}

这一结果就是 \cite{ref54} 的 (1.1) 节中以皮卡群（Picard group）语言陈述的定理。该定理在那里是对多个点 $p_1,\dots,p_s$ 处有抛物结构的情形陈述的。此时存在类似上述的纤维化，在每个点 $p_i$ 处各有一个 $G/\TT$ 的副本，因此类似的分析给出
% 式 (3.48)
\setcounter{equation}{47}
\begin{equation}\label{eq:3.48}
  H^2(\M(\alpha_1,p_1;\dots;\alpha_s,p_s);\Z)=\Z\oplus\left(\oplus_{i=1}^s
\Lambda_{{\rm wt},i}\right),
\end{equation}
其中有 $s$ 个权格的副本。

形如~\eqref{eq:3.47} 的格分裂一般是非典范的，但在本例中存在典范分裂。看出这一点的一种方式，是从伴随表示中的通用丛（universal bundle）$E_{\rm ad}\to \M\times C$ 出发。限制到 $\M\times p$ 上，它给出一个 $G_{\rm ad}$-丛 $E_{{\rm ad},p}\to \M$。拉回到 $\M(\alpha;p)\to \M$ 上时，该丛的结构群约化为 $\TT$，于是它分裂为线丛之和。这些线丛的第一陈类有理地生成~\eqref{eq:3.47} 中的直和项 $\Lambda_{\rm wt}$。这就给出该分裂；在我们下面通过仿射权格给出的解释中，这一点也是显然的。

\begin{yizhu}[易错点]式~\eqref{eq:3.47} 的分裂之所以典范，是因为它由几何直接给出：伴随丛 $E_{\rm ad}\to \M\times C$ 限制到 $\M\times p$ 得 $G_{\rm ad}$-丛 $E_{{\rm ad},p}\to\M$，拉回后结构群约化到 $\TT$，于是按特征分解为线丛之和，这些线丛的第一陈类张成 $\Z\oplus\Lambda_{\rm wt}$ 中的 $\Lambda_{\rm wt}$ 一项。一般的格短正合列（如~\eqref{eq:3.45}）只保证分裂存在，并不指定取法；这里的分裂与后文用仿射权格给出的解释一致，正是因为两者都把 $\Lambda_{\rm wt}$ 实现为具体的陈类。\end{yizhu}

在推导~\eqref{eq:3.47} 时，我们用了近似纤维化~\eqref{eq:3.38}。事实上，我们所忽略的奇点对第二上同调并不重要，除非在使某个消没圈（vanishing cycle）收缩、且 $\M(\alpha;p)$ 的第二贝蒂数变小的那些特殊的 $\alpha$ 值处。关于 $G=SU(2)$ 时出现这一情形的 $\alpha$ 值的描述，见 \cite{ref55}。另一方面，若考虑抛物丛的“叠”（stack）的上同调（而不是稳定抛物丛的模空间 $\M(\alpha;p)$），则不会发生跳变。对规范场论而言这实际上才是正确的对象，因为出发点是所有规范场构成的空间，它是代数几何中所有丛或抛物丛构成的叠在微分几何中的类比物。（在西格玛模型中，支撑在消没圈上的膜在消没圈收缩时并不消失，而只是变成支撑在由此产生的奇点上的膜。）用通俗的话说，$\M_\mathbb{L}(\alpha;p)$ 的第二上同调在某些 $\alpha$ 值处的下降对我们的应用无关紧要，原因如下。当我们描述 $\MH$ 的辛形式 $\omega_I$、$\omega_J$、$\omega_K$ 的上同调类时，跳变只意味着某些周期在某些 $\alpha$ 值处必须为零，这在我们给出的公式中将显而易见。或者，在描述仿射外尔群（affine Weyl group）对 $\MH$ 的上同调的作用时，我们将避开参数的坏值。

\begin{yizhu}[物理图像补充]跳变的机制：$\alpha$ 取特殊值时某个消没圈收缩，$\M(\alpha;p)$ 的第二贝蒂数变小，式~\eqref{eq:3.45} 的中间群随之降秩；若取抛物丛叠的上同调则无此现象，因为叠把收缩前后的对象算作同一族。西格玛模型的图像是：支撑在消没圈上的膜不因圈收缩而消失，只改为支撑在奇点上。$SU(2)$ 情形下跳变发生的 $\alpha$ 值由 Pauly~\cite{ref55} 给出，可用来核对式~\eqref{eq:3.45} 在各连通分支上的适用范围。\end{yizhu}

上述对 $\M(\alpha;p)$ 第二上同调的描述，将使我们能够有效地描述该空间辛形式的上同调类。该空间的辛形式为
% 式 (3.49)
\setcounter{equation}{48}
\begin{equation}\label{eq:3.49}
  \omega=-{1\over 4\pi}\int_C\,\Tr\,\delta A\wedge
\delta A
\end{equation}
它取值于 $H^2(\M(\alpha;p);\R)=H^2(\M(\alpha;p);\Z)\otimes\R=
\R\oplus {\mathfrak t}^\vee$。它的上同调类在 \cite{ref56} 的定理 3.2 与命题 3.7 中给出；关于 $SU(2)$ 情形用粘接论证的证明见 \cite{ref57}，而基于把 $\M(\alpha;p)$ 实现为以 $\alpha$ 为矩映射中参数的辛商的更一般论证，见 \cite{ref58}。我们将在若干准备之后解释该公式。

\bigskip\noindent{\it 仿射外尔群}

我们回忆：在等价意义下，$\alpha$ 取值于 $\TT/{\mathcal W}$，其中 ${\mathcal W}$ 是外尔群。或者，可将 $\alpha$ 提升到 $\mathfrak t$，但此时用余特征格（cocharacter lattice）$\Lambda_{\rm cochar}$ 的元素平移 $\alpha$ 会带来等价。由于已假定 $G$ 单连通，余特征格与余根格（coroot lattice）$\Lambda_{\rm cort}$ 相同。格平移与外尔变换合成的群
% 式 (3.50)
\setcounter{equation}{49}
\begin{equation}\label{eq:3.50}
  \AffWeyl=\Lambda_{\rm cort}\rtimes \Weyl
\end{equation}
称为仿射外尔群。它是 $G$ 的 Kac-Moody 代数或仿射李代数的外尔群。因此可以把 $\alpha$ 看作取值于 $\mathfrak t/\AffWeyl$。

可以如下把 $\mathfrak t$ 划分为 $\AffWeyl$ 作用的基本区域（fundamental domain）。在 $\mathfrak t$ 的某些超平面上，$U=\exp(-2 \pi\alpha)$ 是非正则的（它与 $G$ 的某个非交换子群交换）。每个这样的超平面都是 $\AffWeyl$ 中某个元素的不动点轨迹。这些超平面把 $\mathfrak t$ 分成仿射外尔群作用的基本区域。

例如在 $G=SU(N)$ 的情形，推广~\eqref{eq:3.17}，可把 $\alpha$ 看作对角矩阵 $i\,{\rm diag}(y_1,\dots,y_N)$，其中 $\sum_a y_a=0$。${\mathcal W}$ 通过置换作用，$\Lambda_{\rm cort}$ 通过保持 $\sum_a y_a=0$ 的整数平移作用。$U=\exp(-2\pi\alpha)$ 非正则的条件是：对某对 $a,b$，$n=y_a-y_b$ 为整数；此时 $\alpha$ 在仿射外尔变换 $y_a\to y_b+n$、$y_b\to y_a-n$ 下不变。超平面 $n=y_a-y_b$ 把 $\mathfrak t$ 分成 $\AffWeyl$ 的基本区域。

通常的外尔群在选定一组正根（positive root）之后，在 $\mathfrak t$ 中有一个特殊的基本区域，称为正外尔房（positive Weyl chamber）。它是使 $\langle \alpha,w\rangle >0$ 对每个基本权（fundamental weight）$w$ 都成立的那个区域。对 $G=SU(N)$，该外尔房由 $y_1\geq y_2\geq\dots \geq y_N$ 描述。将其与上述超平面相交，正外尔房便分解为仿射外尔群的无穷多个基本区域之并。其中有一个独特的区域，我们记为 $\mathfrak{D}$，其闭包包含 $\alpha=0$。对 $G=SU(N)$，$\mathfrak{D}$ 由 $y_1-y_N\leq 1$ 刻画（若希望 $U$ 是正则的，则取 $y_1-y_N<1$），这是对~\eqref{eq:3.17} 中条件 $y<1/2$ 的推广。

现在可以解释上面提到的、关于 $\M(\alpha;p)$ 的辛形式 $\omega$ 的上同调类的公式。我们把该公式写在刚才描述的 $\mathfrak t$ 的独特仿射外尔房中。于是上同调类 $[\,\omega/2\pi]$ 为
% 式 (3.51)
\setcounter{equation}{50}
\begin{equation}\label{eq:3.51}
  \left[\,\omega\over 2\pi\right]=e\oplus
(- \alpha^*).
\end{equation}
这里我们用到了 $[\,\omega/2\pi]$ 取值于 $\R\oplus \mathfrak t^\vee$ 这一事实。在~\eqref{eq:3.51} 中，$e\in H^2(\M;\R)\cong\R$ 是行列式线丛 $\mathfrak L\to\M$ 的第一陈类拉回到 $\M(\alpha;p)$ 上的像（等价地，是 $\M$ 的辛形式 $\omega/2\pi$ 的上同调类的拉回）。而右端第二项中，$\alpha^*\in \mathfrak t^\vee$ 是（如同第~\ref{sec:2.4}~节）$\alpha\in \mathfrak t$ 在由二次型 $-\Tr$ 给出的映射 $\mathfrak t\to \mathfrak t^\vee$ 下的像。

\begin{yizhu}[易错点]式~\eqref{eq:3.51} 两项来源不同：$e\in H^2(\M;\R)\cong\R$ 是行列式线丛 $\mathfrak L\to\M$ 的第一陈类（等价地，$\M$ 的辛形式类 $\omega/2\pi$）沿 $\M(\alpha;p)\to\M$ 的拉回；$-\alpha^*\in\mathfrak t^\vee$ 是 $\alpha$ 在二次型 $-\Tr$ 下的像。与式~\eqref{eq:3.53} 对照可见符号一致：那里 $\int_{\Sigma_p}\Tr\alpha F/2\pi=-\langle\alpha^*,\mathfrak m\rangle$，其中 $\mathfrak m\in\Lambda_{\rm cort}$ 是 $F/2\pi$ 的类；把 $\alpha$ 换成 $-\alpha$ 只改第二项的号，第一项不动。这也解释了为何 $e$ 只依赖底流形而 $-\alpha^*$ 携带抛物权重：核对公式时不要把两者混为一项。\end{yizhu}

我们在此不概述 \cite{ref57,ref58} 中导出该公式的论证，而用四维规范场论给出简要说明。为确定 $\omega$ 的上同调类，需计算周期 $\int_\Sigma\omega$，其中 $\Sigma\subset \M(\alpha;p)$ 是闭二维流形。给定 $\Sigma$ 的选取，令 $M$ 为四维流形 $M=\Sigma\times C$，并在 $M$ 上作通常的构造。令 $\Sigma_p=\Sigma \times p$。$\Sigma$ 的每一点确定 $C\backslash p$ 上一个绕 $p\in C$ 有单值性的平坦 $G$-丛；这些丛拼成一个在 $M\backslash \Sigma_p$ 上、绕 $\Sigma_p$ 有单值性的 $G$-丛 $E\to
M\backslash \Sigma_p$。（$E$ 可能只能作为 $G_{\rm ad}$-丛存在，但这不影响下面的推导。）$\M(\alpha;p)$ 的辛形式与四维规范场论之间的关系是
% 式 (3.52)
\setcounter{equation}{51}
\begin{equation}\label{eq:3.52}
  \int_\Sigma{\omega\over 2\pi}=\int_{M}{\Tr\,F\wedge F\over 8\pi^2}
\end{equation}

这一关系式的证明与无抛物结构时的相应陈述完全相同；见 \cite{ref8} 的式 (4.18)。现在利用~\eqref{eq:2.41}（并注意到 $\Sigma_p\cap \Sigma_p=0$），我们有
% 式 (3.53)
\setcounter{equation}{52}
\begin{equation}\label{eq:3.53}
  \int_{\Sigma}{\omega\over 2\pi}=\int_{\Sigma_p}\Tr\,{\alpha}{F\over 2\pi}
     =-\langle \alpha^*,\mathfrak{m}\rangle~{\rm mod}~\Z,
\end{equation}
其中 $\mathfrak{m}\in H^2(\Sigma_p,\pi_1(\TT))\cong \Lambda_{\cort}$ 是 $F/2\pi$ 的上同调类，$\langle~,~\rangle$ 是 $\mathfrak t$ 与 $\mathfrak t^\vee$ 之间的自然配对。这与~\eqref{eq:3.51} 等价。我们还看到，能与辛形式自然配对的同调圈（homology cycle）由 $\Lambda_{\cort}$（$\Lambda_{\wt}$ 的对偶）标记。

对 $\M(\alpha;p)$ 辛形式的这一描述，与我们在第~\ref{sec:3.2}~节关于希格斯丛的论断有明显的相似性。对希格斯丛，我们预期 $[\,\omega_I]$ 之类的上同调类线性变化，如~\eqref{eq:3.8} 所断言；而对普通丛，我们得到~\eqref{eq:3.51} 中看到的线性性。我们将在第~\ref{sec:3.6}~节解释这些结果之间的关系。

然而，在希格斯丛情形，导出~\eqref{eq:3.8} 的论证对所有 $\alpha$ 成立；而丛情形的类似陈述~\eqref{eq:3.51} 只对位于某个基本仿射外尔房中的 $\alpha$ 成立。对丛而言，没有明显的方式把公式~\eqref{eq:3.51} 延拓到基本仿射外尔房之外，因为在 $\mathfrak{D}$ 的边界上，流形 $\M(\alpha;p)$ 坍缩为更低维的流形，正如我们将在第~\ref{sec:3.5}~节看到的。对希格斯丛，如我们将在第~\ref{sec:3.6}~节与第~\ref{sec:3.7}~节描述的，可取 $\beta,\gamma\not=0$，从而光滑地延拓到仿射外尔房边界之外。

\begin{yizhu}[含糊点拨]“正如我们将在第~\ref{sec:3.5}~节看到的”一句是本节内部的前向引用：坍缩为低维流形的具体讨论在本片段稍后（原文 PDF 第 61 页，印刷第 60 页，$\alpha\to\bar\alpha$ 一段）才出现。作者在此提前引用，读者不必到别节寻找。另外，“微分消失”的完整论证也要到“仿射权格”一段（PDF 第 57--58 页）才补出，那里用 $H^2(\M(\alpha;p);\Z)$ 必须含仿射权格来反推 $d_2=0$。\end{yizhu}

在上述讨论中，我们暗中取列维子群为 $\mathbb{L}=\TT$。然而同一结果对任意 $\mathbb{L}$ 成立，推导相同；只需把 $\alpha$ 限制为 $\mathbb{L}$-不变的。

\bigskip\noindent{\it 仿射权格}

重新考虑 $\mathbb{L}=\TT$ 情形下 $\M(\alpha;p)$ 第二上同调的描述~\eqref{eq:3.47}：
% 式 (3.54)
\setcounter{equation}{53}
\begin{equation}\label{eq:3.54}
  H^2(\M(\alpha;p);\Z)=\Z\oplus
\Lambda_{wt}.
\end{equation}
我们可以这样描述：$H^2(\M(\alpha;p);\Z)$ 是 $G$ 的仿射权格（affine weight lattice），即与 $G$ 相伴的仿射李代数、Kac-Moody 代数或中心扩张圈群的权格。这一描述在现有数学理论中是自然的~\cite{ref54}；在那里~\eqref{eq:3.54} 由圈群理论得到。我们将在 $2+1$ 维 Chern-Simons 规范场论的语境中，用物理语言勾勒这一想法~\cite{ref59}。令 $k$ 为正整数，$w\in \Lambda_{wt}$ 为 $G$ 的 Kac-Moody 代数在能级 $k$ 的可积表示的最高权。用二次型 $-\Tr$ 把 $\mathfrak t^\vee={}^L\mathfrak t$ 与 $\mathfrak t$ 等同，$w\in {}^L\mathfrak t$ 映到元素 $w^*\in \mathfrak t$。考虑 $C$ 上的 Chern-Simons 规范场论，带一个标记点 $p$，标以 $G$ 的最高权为 $w$ 的表示。经典相空间是取 $\alpha=w^*/k$ 的 $\M(\alpha;p)$。量子化通过取线丛 $\mathfrak L^k$ 的整体截面来实现，其第一陈类在 de Rham 上同调中由 $k\omega/2\pi$ 代表，其中 $\omega=-{1\over 4\pi}\int_C\Tr\,\delta A\wedge \delta A$。（除非 $w=0$，$\mathfrak L^k$ 并不真的是某个线丛的 $k$ 次幂，但由于有理第一陈类中含有因子 $k$，我们仍把它写成 $\mathfrak L^k$。）因此，de Rham 上同调中的类 $k\omega/2\pi$ 提升为格 $H^2(\M(\alpha;p);\Z)$ 中的一个元素。相应地，该格必须包含对应于对 $(k,w)$ 的点，也就是说，它必须包含 $G$ 的仿射权格。由关于 $[\,\omega/2\pi]$ 的上同调类的公式~\eqref{eq:3.51} 可知，这些点在 de Rham 上同调中的像互不相同。至此也很清楚：关于 $H^2(\M(\alpha;p);\Z)$ 的谱序列~\eqref{eq:3.40} 中的微分确实消失（正如我们在上述讨论中所假定的），否则该上同调群不可能包含仿射权格。

\begin{yizhu}[物理图像补充]这一段其实给出了 $d_2=0$ 的物理反推：对每个能级 $k$ 与最高权 $w$，经典相空间是 $\alpha=w^*/k$ 的 $\M(\alpha;p)$，量子化要求 $k\omega/2\pi$ 为整类（即线丛 $\mathfrak L^k$ 的第一陈类）；由式~\eqref{eq:3.51}，它在 de Rham 上同调中的像是 $k e\oplus(-w^*)$。不同的 $(k,w)$ 给出无穷多个互不相同的点，故 $H^2(\M(\alpha;p);\Z)$ 必须容纳整个仿射权格；若谱序列的 $d_2$ 非零，$H^2$ 的秩将小于 $\Z\oplus\Lambda_{\rm wt}$，容纳不下。注意 $w=0$ 时 $\mathfrak L^k$ 才真是某个线丛的 $k$ 次幂。\end{yizhu}

由于 $H^2(\M(\alpha;p);\Z)$ 是 $G$ 的仿射权格，它承载 $G$ 的仿射外尔群的自然作用，尽管该作用在 $\M(\alpha;p)$ 的几何中没有显然的意义。我们将在第~\ref{sec:3.7}~节更概念化地解释 $H^2(\M(\alpha,p);\Z)$ 为何承载仿射外尔群的这一作用，并将更精确地描述该作用。

若放松 $G$ 单连通的假定，会发生什么？通过 Kac-Moody 代数的论证表明：对任意 $G$（不必单连通），格 $H^2(\M(\alpha,p;G); \Z)$ 模去挠（torsion）后包含 $G$ 的仿射特征格（affine character lattice）作为子格。所谓单的（但未必单连通的）群 $G$ 的仿射特征格，我们指的是
% 式 (3.55)
\setcounter{equation}{54}
\begin{equation}\label{eq:3.55}
  \Lambda_{\rm aff\,char}=\Z\oplus
\Lambda_{\rm char}(G),
\end{equation}
其中 $\Z$ 分类 $G$ 的圈群（loop group）的中心扩张（central extension）\footnote{此处需稍加小心。令 $\bar G$ 为 $G$ 的泛覆盖。由于瞬子数对 $\bar G$ 取 $\Z$ 值而对 $G$ 不然，并非 $\bar G$ 的圈群的每个中心扩张都对应 $G$ 的一个中心扩张。\eqref{eq:3.55} 中的直和项 $\Z$ 自然地理解为~\eqref{eq:3.54} 中直和项 $\Z$ 的真子群。}。我们猜测，对通有的 $\alpha$，$H^2(\M(\alpha,p;G);\Z)$ 模去挠后总是恰为该格：
% 式 (3.56)
\setcounter{equation}{55}
\begin{equation}\label{eq:3.56}
  H^2(\M(\alpha;p))=\Z\oplus \Lambda_{\rm char}(G).
\end{equation}

\begin{yizhu}[易错点]非单连通时只有包含关系，且原文明确标为猜测。其一，$\M(\alpha,p;G)$ 的上同调含 $\pi_1(G)$ 带来的挠部，\eqref{eq:3.55} 说的是模去挠之后的自由部分，不能与格直接比较。其二，\eqref{eq:3.55} 中的 $\Z$ 并不是~\eqref{eq:3.54} 中的那个 $\Z$：圈群的每个中心扩张未必下降为 $G$ 的中心扩张，因为在泛覆盖 $\bar G$ 上瞬子数取 $\Z$ 值而在 $G$ 上不然，脚注给出的正是这个理由。所以式~\eqref{eq:3.56} 只是猜测，引用时不可当作定理。\end{yizhu}

无论如何，我们将在第~\ref{sec:3.7}~节证明：该格的一个自然对称群作用在 $H^2(\M(\alpha,p;G);\Z)$ 上。结果表明：若限制在 $\M(\alpha,p;G)$ 的一个连通分支上，则无论 $G$ 是否单连通，作用在上同调上的群都是同一个仿射外尔群。

\bigskip\noindent{\it 另一视角}

近似纤维化~\eqref{eq:3.39} 也可以从纯拓扑或辛几何的观点看出，而不必提及抛物结构（例如见 \cite{ref56} 中的简要概述）。在黎曼面（Riemann surface）$C$ 上取一组标准的 $A$-圈与 $B$-圈。令 $A_i$、$B_i$ 为平坦 $G$-丛绕这些圈的单值性。它们满足
% 式 (3.57)
\setcounter{equation}{56}
\begin{equation}\label{eq:3.57}
  \prod_{i=1}^g[A_i,B_i]=1,
\end{equation}
其中 $[A,B]=ABA^{-1}B^{-1}$。平坦丛的模空间 $\M$ 是该方程的解空间模去共轭。取 $G$ 中接近单位元的元素 $U$，把~\eqref{eq:3.57} 换成
% 式 (3.58)
\setcounter{equation}{57}
\begin{equation}\label{eq:3.58}
  \prod_{i=1}^g[A_i,B_i]=U.
\end{equation}
该方程的解描述 $C\backslash p$ 上一个绕 $p$ 有单值性 $U$ 的平坦丛。只要不遇到 $\M$ 的奇点，当 $U$ 足够接近 $1$ 时，\eqref{eq:3.57} 与~\eqref{eq:3.58} 的解空间在拓扑上相同。现在假设我们只想指定 $U$ 的共轭类，并假定它是半单的，且同样足够接近 $1$。那么 $U$ 的共轭类包含元素 $\exp(-2\pi\alpha)$，其中 $\alpha$ 是 $\mathfrak t$ 中接近原点的元素。令 $\mathbb{L}$ 为 $G$ 中与 $\alpha$（等价地，与 $U$）交换的子群；则 $\mathbb{L}$ 是列维子群。$U$ 在其共轭类中的可能选取由一个 $G/\mathbb{L}$ 的副本参数化，因此~\eqref{eq:3.58} 的解空间是~\eqref{eq:3.57} 的解空间上的 $G/\mathbb{L}$-丛。对两个空间都模去 $G$ 的共轭作用后依然如此，于是我们再次得到纤维化~\eqref{eq:3.39}：
% 式 (3.59)
\setcounter{equation}{58}
\begin{equation}\label{eq:3.59}
  \begin{matrix}G/\mathbb{L}&\to& \M_\mathbb{L}(\alpha;p)\\
                                &   & \downarrow \\
                                & & ~\M. \end{matrix}
\end{equation}

假设我们要从 $C$ 中挖去点 $p_1,\dots,p_s$，并记 $U_1,\dots,U_s$ 为绕这些点的单值性。\eqref{eq:3.57} 的类比是
% 式 (3.60)
\setcounter{equation}{59}
\begin{equation}\label{eq:3.60}
  \prod_{i=1}^g[A_i,B_i]=\prod_{a=1}^s U_a.
\end{equation}
若各 $U_a$ 足够接近单位元且与列维子群 $\mathbb{L}_a$ 交换，则类似上述的论证给出一个在每个穿刺点处带一个 $G/\mathbb{L}_a$ 因子的纤维化（例如这导出了式~\eqref{eq:3.48} 中 $\M(\alpha_1,p_1;\dots;\alpha_s,p_s)$ 第二上同调的描述）。然而，我们想考虑另一个面向后续应用的问题。令 $y_a$（$a=1,\dots,s$）取值于 $G$ 的中心，并假定
% 式 (3.61)
\setcounter{equation}{60}
\begin{equation}\label{eq:3.61}
  \prod_{a=1}^sy_a=1.
\end{equation}
则式~\eqref{eq:3.60} 在
% 式 (3.62)
\setcounter{equation}{61}
\begin{equation}\label{eq:3.62}
  U_a\to y_aU_a,~~a=1,\dots,s.
\end{equation}
下完全不变。

若 $U_a$ 共轭于某个 $\alpha_a\in \mathfrak t$ 的 $\exp(-2\pi \alpha_a)$，且 $y_a=\exp(-2\pi u_a)$，其中 $u_a\in \Lambda_{\rm cowt}$，则变换~\eqref{eq:3.62} 相当于
% 式 (3.63)
\setcounter{equation}{62}
\begin{equation}\label{eq:3.63}
  \alpha_a\to \alpha_a+u_a.
\end{equation}
而~\eqref{eq:3.61} 等价于
% 式 (3.64)
\setcounter{equation}{63}
\begin{equation}\label{eq:3.64}
  \sum_a u_a\in \Lambda_{\rm cochar}\subset \Lambda_{\rm cowt}.
\end{equation}
这是式~\eqref{eq:2.6} 与~\eqref{eq:2.7} 所描述情形的一个例证。各个平移 $\alpha_a\to \alpha_a+u_a$ 单独会改变由 $E$ 导出的 $G_{\rm ad}$-丛的特征类 $\xi$，而条件~\eqref{eq:3.61} 或~\eqref{eq:3.64} 保证 $\xi$ 整体上实际不变。

\begin{yizhu}[含糊点拨]式~\eqref{eq:3.61}、\eqref{eq:3.64} 是第~\ref{sec:2.1}~节式~\eqref{eq:2.6}、\eqref{eq:2.7} 的多穿刺点版本。那里单点时有 $\xi\to\xi+y[D]$，多个 $D_i$ 时要求 $\sum_iy_i[D_i]=0$；本节把 $D_i$ 换成 $\Sigma_{p_a}=\Sigma\times p_a$，各 $[\Sigma_{p_a}]$ 线性无关，于是条件化为 $\prod_a y_a=1$ 与 $\sum_a u_a\in\Lambda_{\rm cochar}$。符号约定略有不同：第~\ref{sec:2.1}~节写 $y_i=\exp(2\pi u_i)$，本节因取单值性 $\exp(-2\pi\alpha_a)$ 而写 $y_a=\exp(-2\pi u_a)$，差一符号但不影响结论。这一判据的物理含义是：单个穿刺点的平移各自独立，但若要在同一个 $G$-丛上同时实现全部平移，特征类的总变化必须为零，否则 $G_{\rm ad}$-丛无法整体定义。\end{yizhu}

现在我们为简单起见回到单个穿刺点的情形。若 $U=\exp(-2\pi\alpha)$ 是正则的（这可把 $\alpha$ 放在独特仿射外尔房 $\mathfrak{D}$ 的内部来实现），则纤维化~\eqref{eq:3.59} 取形式
% 式 (3.65)
\setcounter{equation}{64}
\begin{equation}\label{eq:3.65}
  \begin{matrix}G/\TT&\to& \M(\alpha;p)\\
                                &   & \downarrow \\
                                & & ~\M. \end{matrix}
\end{equation}

现在假设 $\alpha$ 趋近 $\mathfrak{D}$ 的一个边界点 $\bar\alpha$。此时 $U$ 不再正则，其共轭轨道是对某个 $\mathbb{L}$ 而言的 $G/\mathbb{L}$ 副本，而不再是 $G/\TT$。于是纤维化现在取~\eqref{eq:3.59} 中的形式：
% 式 (3.66)
\setcounter{equation}{65}
\begin{equation}\label{eq:3.66}
  \begin{matrix}G/\mathbb{L}&\to& \M_\mathbb{L}(\bar\alpha;p)\\
                                &   & \downarrow \\
                                & & ~\M. \end{matrix}
\end{equation}

当 $\alpha\to\bar\alpha$ 时会发生什么？$G/\TT$ 纤维化于 $G/\mathbb{L}$ 之上，纤维为 $\mathbb{L}/\TT$。当 $\alpha\to\bar\alpha$ 时，$U$ 的轨道从 $G/\TT$ 的一个副本坍缩为 $G/\mathbb{L}$ 的一个副本，于是我们得到纤维化
% 式 (3.67)
\setcounter{equation}{66}
\begin{equation}\label{eq:3.67}
  \begin{matrix}\mathbb{L}/\TT&\to& \M(\alpha;p)\\
                                &   & \downarrow \\
                                & & ~\M_\mathbb{L}(\bar\alpha;p). \end{matrix}
\end{equation}
纤维 $\mathbb{L}/\TT$ 是辛流形，它们在 $\alpha\to\bar\alpha$ 时收缩为点。因此，当 $\alpha$ 趋近非正则值 $\bar\alpha$ 时，$\M(\alpha;p)$ 发生的就是这种事：它坍缩为更低维的簇（variety），其消没圈为 $\mathbb{L}/\TT$。

在通有的情形下，边界点 $\bar\alpha$ 只包含在标记 $\mathfrak{D}$ 边界的某一个超平面中，此时 $\mathbb{L}/\TT=SU(2)/U(1)=\mathbb{CP}^1=S^2$。例如，对 $G=SU(N)$，有 $\alpha=i\,{\rm diag}(y_1,y_2,\dots,y_N)$，其中 $y_1\geq y_2\geq y_3\dots,\geq y_N$、$y_1-y_N\leq 1$、$\sum y_a=1$。在通有的边界点处，恰有一个不等式取等号，此时 $\mathbb{L}=U(2)\times U(1)^{N-3}$，且 $\mathbb{L}/\TT=U(2)/U(1)^2=\mathbb{CP}^1=S^2$。因此当 $\alpha$ 趋近 $\mathfrak{D}$ 的通有边界点时，$\M(\alpha;p)$ 以消没的二维球面（即收缩为点的二维球面）为纤维。一般地，$\mathbb{L}$ 的半单部分的秩等于包含 $\bar\alpha$ 的边界超平面的个数。

\begin{yizhu}[原文疑误]此处写 $\sum y_a=1$，但本片段前面（原文 PDF 第 56 页，印刷第 55 页）以及 $SU(N)$ 的约定均为 $\sum_a y_a=0$：$\alpha=i\,{\rm diag}(y_1,\dots,y_N)$ 须无迹。由于只有差 $y_a-y_b$ 有意义，$\sum y_a$ 的归一化不影响后续结论（$y_1-y_N\leq 1$ 与 $\mathbb{L}$ 的判定都只依赖差），但前后不一致仍应指出。原文如此（PDF 第 61 页，印刷第 60 页），照录。\end{yizhu}

\begin{yizhu}[物理图像补充]消没圈的几何：$\bar\alpha$ 落在 $\mathfrak D$ 的边界超平面上时，中心化子由 $\TT$ 涨到更大的 $\mathbb{L}$，$G/\TT\to G/\mathbb{L}$ 的纤维 $\mathbb{L}/\TT$ 在 $\alpha\to\bar\alpha$ 时缩为一点，这就是式~\eqref{eq:3.67} 的消没圈。若 $\bar\alpha$ 通有地只落在一个超平面上，则 $\mathbb{L}/\TT=SU(2)/U(1)=\mathbb{CP}^1$；若 $\bar\alpha$ 同时落在 $j$ 个超平面上，则 $\mathbb{L}$ 的半单部分秩为 $j$，消没圈的维数随之增大。$SU(N)$ 的主导情形即 $\mathbb{L}=U(2)\times U(1)^{N-3}$、$\mathbb{L}/\TT=\mathbb{CP}^1$，与文中“恰有一个不等式取等号”对应。\end{yizhu}

更一般地，可以考虑一对列维子群 $\mathbb{L}_1$ 与 $\mathbb{L}_2$，其中 $\mathbb{L}_1$ 是 $\mathbb{L}_2$ 的真子群，但不一定等于 $\TT$。与上述同类的论证仍然适用。若 $U=\exp(-2\pi\alpha_1)$，其中 $\alpha_1$ 是 $\mathbb{L}_1$-正则的，则得到纤维化
% 式 (3.68)
\setcounter{equation}{67}
\begin{equation}\label{eq:3.68}
  \begin{matrix}G/\mathbb{L}_1&\to& \M_{\mathbb{L}_1}(\alpha_1;p)\\
                                &   & \downarrow \\
                                & & ~\M. \end{matrix}
\end{equation}
若 $\alpha_2$ 是 $\mathbb{L}_2$-正则的，则对 $U=\exp(-2\pi\alpha_2)$，纤维化形如
% 式 (3.69)
\setcounter{equation}{68}
\begin{equation}\label{eq:3.69}
  \begin{matrix}G/\mathbb{L}_2&\to& \M_{\mathbb{L}_2}(\alpha_2;p)\\
                                &   & \downarrow \\
                                & & ~\M. \end{matrix}
\end{equation}
当 $\alpha_1$ 沿 $\mathbb{L}_1$-正则值趋近 $\alpha_2$ 时，$U$ 的轨道 $G/\mathbb{L}_1$ 退化为 $G/\mathbb{L}_2$。$G/\mathbb{L}_1$ 以 $\mathbb{L}_2/\mathbb{L}_1$ 为纤维映到 $G/\mathbb{L}_2$；而当 $\alpha_1$ 趋近 $\alpha_2$ 时，在奇点之外我们得到纤维化
% 式 (3.70)
\setcounter{equation}{69}
\begin{equation}\label{eq:3.70}
  \begin{matrix}\mathbb{L}_2/\mathbb{L}_1&\to& \M_{\mathbb{L}_1}(\alpha_1;p)\\
                                &   & \downarrow \\
                                & & ~\M_{\mathbb{L}_2}(\alpha_2;p). \end{matrix}
\end{equation}
纤维 $\mathbb{L}_2/\mathbb{L}_1$ 是 $\alpha_1\to\alpha_2$ 的消没圈。

正如所料，极限行为也有一个用复几何给出的自然描述。由 $(\mathbb{L}_1,\alpha_1)$ 决定一个抛物子群 $\EUP_1$，由 $(\mathbb{L}_2,\alpha_2)$ 决定一个包含 $\EUP_1$ 的抛物子群 $\EUP_2$。{}从全纯的观点看，根据 Mehta 与 Seshadri 的定理，$\M_{\mathbb{L}_1}(\alpha_1;p)$ 或 $\M_{\mathbb{L}_2}(\alpha_2;p)$ 分别参数化在 $p$ 处带有 $\EUP_1$ 型或 $\EUP_2$ 型抛物结构的丛。把丛的结构群约化到 $\EUP_2$ 之后，进一步约化到子群 $\EUP_1$ 的方式由 $\EUP_2/\EUP_1$ 参数化。于是从全纯的观点看，我们得到纤维化
% 式 (3.71)
\setcounter{equation}{70}
\begin{equation}\label{eq:3.71}
  \begin{matrix}\EUP_2/\EUP_1&\to& \M_{\mathbb{L}_1}(\alpha_1;p)\\
                                &   & \downarrow \\
                                & & \M_{\mathbb{L}_2}(\alpha_2;p). \end{matrix}
\end{equation}
这与~\eqref{eq:3.70} 相符，因为 $\EUP_2/\EUP_1=\mathbb{L}_2/\mathbb{L}_1$。该纤维化即 \cite{ref60} 中的命题 3.4。

\begin{yizhu}[易错点]式~\eqref{eq:3.70} 与~\eqref{eq:3.71} 是同一纤维化的两种描述，相容性来自 $\EUP_2/\EUP_1\cong\mathbb{L}_2/\mathbb{L}_1$：因 $\EUP_2=\mathbb{L}_2\ltimes\mathfrak u_2$ 且 $\mathfrak u_2\subset\EUP_1$（$\EUP_1\subset\EUP_2$），故 $\EUP_2/\EUP_1\cong\mathbb{L}_2/(\EUP_1\cap\mathbb{L}_2)=\mathbb{L}_2/\mathbb{L}_1$。这正是 Boden--Hu~\cite{ref60} 命题 3.4 的内容：$\alpha_1\to\alpha_2$ 的消没圈在复、辛两种描述下都是 $\mathbb{L}_2/\mathbb{L}_1$。原文把这一步视为显然；注意嵌套 $\EUP_1\subset\EUP_2$ 是必要前提，只有此时商 $\EUP_2/\EUP_1$ 才有意义，也才能把 $\alpha_1\to\alpha_2$ 的极限写成两族模空间之间的纤维化。\end{yizhu}
